\section{A Joint Evaluation of Executed Cover}\label{sec:experiment}
\textbf{Cover generation changes residual inference and resource use, while business correctness requires a separate endpoint.}
The five-arm live study compares dispatch policies on actual TLS records. A record-only extension withdraws call boundaries; a separate earlier state control tests whether completion alone can miss collateral effects. These populations and their statistical roles remain separate.
\subsection{Actual dispatch and timing-aware attacks}\label{sec:cover-study}
\textbf{Immediate controller cover is removable because it can skip the model wait preceding honest calls.}
Two fixed-weight planners, Qwen3:8B and Llama3.1:8B, each execute 300 fixtures in five conditions, yielding 3,000 main episodes and 600 additional shadow chains. The six business actions use seven MCP services through LangGraph over TLS~1.3/HTTP/2. Original partitions contain 120 training, 60 validation and 120 paired test fixtures per cell. The task label is requested intent; failures remain included. Cover is isolated from real business state and hidden from the real planner.

No cover establishes the baseline. Immediate cover dispatches after an honest response. Historical sequential and concurrent policies sample training-only honest gaps and service durations, wait in real time and hold responses for the sampled minimum duration. The concurrent dispatcher overlaps real model inference. Shadow runs a complete second planner on an independent six-class ticket; either chain can start first. Calls to the same service serialize. The original observer receives destinations, ciphertext bytes, TLS record completion times and supplied episode/call windows. Section~\ref{sec:no-call} removes the last of these call-level privileges.

\begin{table*}[t]\centering\small
\caption{Joint evaluation of the five live policies, 120 test fixtures per row; uniform six-class reference 20/120. Sel. is the validation-selected non-timing/adaptive count; Max is the largest test count in the frozen suite, reported descriptively, not substituted into the primary tests. GMM is role-label-free removal. Bytes gives whole-connection/call-window factors; Tokens includes both planners. Goal and State are separate counts. In shadow, both selectors choose the identical non-timing attack.}\label{tab:live-timing-main}
\setlength{\tabcolsep}{4pt}
\begin{tabularx}{\textwidth}{@{}lYrrrrrr@{}}\toprule
Planner & Policy & Sel. S/A & Max & GMM & Bytes$\times$ & Tokens$\times$ & Goal/State \\\midrule

Qwen3:8B & No cover & 117/117 & 117 & -- & 1.00/1.00 & 1.00 & 113/114 \\
Qwen3:8B & Immediate & 75/112 & 112 & 112 & 1.15/1.33 & 1.00 & 111/112 \\
Qwen3:8B & History, sequential & 76/93 & 99 & 53 & 1.15/1.32 & 1.00 & 112/113 \\
Qwen3:8B & History, concurrent & 66/99 & 99 & 38 & 1.15/1.33 & 1.00 & 111/112 \\
Qwen3:8B & Planner shadow & 54/54 & 70 & 34 & 1.31/1.68 & 2.01 & 112/113 \\
Llama3.1:8B & No cover & 110/110 & 112 & -- & 1.00/1.00 & 1.00 & 82/70 \\
Llama3.1:8B & Immediate & 80/108 & 108 & 108 & 1.13/1.27 & 1.01 & 83/67 \\
Llama3.1:8B & History, sequential & 82/104 & 104 & 14 & 1.13/1.26 & 1.01 & 84/67 \\
Llama3.1:8B & History, concurrent & 68/106 & 106 & 66 & 1.13/1.27 & 1.02 & 82/67 \\
Llama3.1:8B & Planner shadow & 64/64 & 64 & 43 & 1.35/1.71 & 1.99 & 82/67 \\
\bottomrule\end{tabularx}\end{table*}


\paragraph{Removal and adaptation}
An unlabelled training-gap mixture selects the short-gap component for removal; a no-cover-trained task classifier decodes the remaining trace. In the immediate arm this restores 112/120 and 108/120 correct predictions, versus selected non-timing references 75 and 80. Both meet the prespecified replication criterion and pass Holm correction. Independent TCP reassembly matches all 773,450 recorded TLS record types, lengths and counts. Recomputing the GMM attack with packet-capture clocks reproduces the scores; software completion timestamps retain measured error (Supplement S2).

Historical sampling changes the short-gap attack, but adapted context removal or joint decoding still recognizes 93/120 and 104/120 sequential executions. The Qwen comparison with its 76/120 non-timing reference has Holm $p=.0559$; Llama's comparison with 82/120 has $p=.00078$. Concurrent adapted counts are 99/120 and 106/120, against 66 and 68. Their 27.5- and 31.7-point \emph{attacker advantages under the same policy} both pass correction. The suite maxima in Table~\ref{tab:live-timing-main} additionally expose Qwen's sequential joint decoder at 99/120, rather than replace the frozen selection after testing.

The policy samples minimum waits, not achieved timing distributions. Qwen's mean 40.5\,ms sampled hold becomes 100.4\,ms sequential and 171.2\,ms concurrent at the record observer. A failed simple mixture is therefore not evidence of successful concealment. In the saved packet-clock diagnostic, Llama's sequential mixture deletes 477/634 honest calls while detecting only 64/182 cover calls, explaining its 14/120 task count. Adapted joint decoding on the same condition reaches 104/120.

\paragraph{Planner-generated cover and cost}
In the call-window suite, shadow's strongest attack is the size/destination RF: 70/120 Qwen and 64/120 Llama, or 58.3\% and 53.3\%, versus the 16.7\% uniform reference. A chain-separation attack scores 49 and 34 and exactly recovers the real chain in 29 and 32 fixtures. The validation selectors choose the same non-timing attack for both roles; their equality is not an independent role-indistinguishability test. Shadow uses 2.01/1.99 times model tokens, 1.68/1.71 times call-window bytes and mean real-chain times 4.80/3.94\,s versus baseline 2.91/3.33\,s. It is a measured operating point, not an equal-budget demonstration that planning itself is superior. Supplement S2 retains every attack, primary test and cost.

\subsection{Recognition without supplied call boundaries}\label{sec:no-call}
\textbf{Withdrawing call endpoints tests a different attacker representation, while whole-episode association remains available.}
We reanalyze all 3,000 recorded episodes using complete encrypted connections, including initialization and closing records. The attacker receives only relative record times, directions, ciphertext lengths and destination identities. It receives no call starts/ends, operation names, role labels or HTTP stream assignment. A per-connection idle-gap rule proposes events; five gaps from 5 to 250\,ms compete on validation task accuracy. Raw-record classifiers bypass this eventization entirely.

For each condition, LR and RF use destination summaries and early/late event sequences, with and without timing. A training-only GMM removes short-gap events before a classifier trained on no-cover traces. Validation separately chooses the baseline and removal attacks. This is a retrospective, fixed-protocol analysis of reused fixtures; only the immediate-arm paired comparison is a primary test, with Holm over the two models. Table~\ref{tab:no-call-main} also gives each finite suite's largest test count without substituting that maximum into the primary test.

\begin{table}[t]\centering\small
\caption{No supplied call boundaries. All encrypted records in associated episode connections enter the analysis. Cells give selected idle-grid baseline / removal / largest test count including the later exploratory protocol-aware suite, out of 120; uniform reference20. Whole-task association remains supplied.}\label{tab:no-call-main}
\begin{tabularx}{\columnwidth}{@{}Yrr@{}}\toprule
Policy & Qwen & Llama \\\midrule
No cover & 115/109/117 & 108/101/112 \\
Immediate & 80/69/91 & 90/86/97 \\
History, sequential & 73/74/75 & 88/56/93 \\
History, concurrent & 56/61/62 & 78/62/88 \\
Planner shadow & 69/65/73 & 69/56/69 \\
\bottomrule\end{tabularx}\end{table}

On immediate cover, selected removal recognizes 69/120 Qwen and 86/120 Llama tasks, against baselines 80 and 90. The paired differences are -9.2 and -3.3 points (Holm $p=0.02545$ and $0.5413$). Removal underperforms the baseline in Qwen after correction; Llama does not reject. The original removal advantage is not reproduced by this boundary-free idle-gap pipeline. After observing the idle-grid results, an exploratory protocol-aware eventizer uses the public request-padding quantum and MCP setup sequence: large upward records mark proposed requests, and large downward records delimit responses before the next request. It does not read call truth. Within this representation, selected removal changes Qwen recognition from 68 to 91, but Llama from 90 to 87; these exploratory counts do not replace the original tests. Immediate cover permits maxima 91/120 and 97/120 across the record-only suites. Whole-episode association remains given, and differing representations prevent interpreting the change from the earlier call-window count as a pure causal boundary effect. Supplement S3 gives the full protocol and independent event-matching diagnostic.

The independent event audit gives protocol-event F1 0.295 and 0.380 at IoU .5 on immediate cover. Task recognition from record summaries therefore does not imply accurate recovery of calls.

\subsection{Business constraints and the protection decision}
Every live cover arm executes in isolated state and hides cover feedback from the real planner. Table~\ref{tab:live-timing-main} reports goal/state counts on the same fixtures as inference and costs. Shadow outcomes are 112/113 for Qwen and 82/67 for Llama, versus no-cover 113/114 and 82/70. These are different endpoint counts: a correctly preserved state need not contain a successful intended action. Low Llama baseline performance and changed planning paths limit business-equivalence conclusions.

The earlier six-policy study supplies a deliberately adverse state control (Supplement S10). Read/proxy/hidden-isolated cover retains Qwen state correctness at 115--116/120. Live writes can alter other objects and reduce it to 7/120; 95 of 102 completed Qwen tasks have collateral errors. Llama's corresponding counts are 3/120 intact and 73 of 76 completed tasks with collateral errors. The harm direction is built into that stress condition; its role is to test whether the evaluation records forbidden effects, not estimate production harms or reproduce WRA decoys.

The combined decision is concrete. Immediate cover leaves high inference after removal. Historical dispatch changes that attack but retains substantial adapted inference. Shadow execution trades additional model work for a different observed trace, with residual task recognition. An adverse write can fail the state contract even when the requested goal passes. Choosing a policy requires stating acceptable inference, task, state and resource tolerances rather than comparing a single accuracy decrease.
